\begin{table}[!t]
\caption{Practical deployment checks for operational PV forecasting.}
\label{tab:deployment_checklist}
\centering
\scriptsize
\setlength{\tabcolsep}{3pt}
\begin{tabularx}{\columnwidth}{@{}>{\raggedright\arraybackslash}p{0.24\columnwidth}>{\raggedright\arraybackslash}X@{}}
\toprule
Layer & Required deployment check \\
\midrule
Data ingest & Verify timezone, cadence, duplicate handling, and missing-value policy before feature generation. \\
Weather input & Record whether target-time weather comes from measured data, API forecasts, numerical weather prediction, or scenarios. \\
Cloud layers & Confirm low-, mid-, and high-cloud availability or document the fallback used to synthesize missing layers. \\
Solar features & Keep latitude, timezone, solar-time convention, and clear-sky proxy consistent between training and inference. \\
Model files & Store Stage-1, Stage-2, sensor, metadata, and feature-list files together to preserve feature order. \\
Post-processing & Apply the same clipping rule used in Eq.~\eqref{eq:postprocess} and state whether metrics are pre- or post-clipping. \\
Monitoring & Track rolling MAE/RMSE, bias, ramp-event errors, and regime-specific errors after deployment. \\
Retraining & Recalibrate after seasonal drift, weather-source changes, plant changes, or sensor replacement. \\
Reporting & Report previous-day persistence, last-value persistence, and sensor-assisted diagnostics separately. \\
\bottomrule
\end{tabularx}
\end{table}
